\section{Proofs of Resource-Optimization Results}
\label{app:optimization-proofs}

This appendix supplies the finite arguments behind the resource results in
Section~4.  All comparisons hold at the fixed belief, primitive model, and
eligible catalog displayed with the corresponding result.

\subsection{Minimum-resource representation and local/global boundary}
\label{app:safe-optimization-proof}

\begin{proof}[Proof of \cref{thm:minimum-safe-compression}.]
The source \(L\) belongs to \(\mathcal P_{\mathrm{safe}}(L)\), so the feasible
family is nonempty.  It is a subset of the finite powerset of \(L\), hence the
rational burden attains a minimum.  Every feasible \(L'\) satisfies
\(K_{L'}=K_L\), so the forward projection theorem identifies all productive
finite-horizon values.  Bellman intertwining and discounted contraction give
the stationary value and complete optimal-action correspondence.
\end{proof}

\begin{proof}[Proof of \cref{thm:local-global-safe-compression}, forward
implication.]
If a global minimum retained a safely deletable active strategy \(s\), then
positivity of \(w_s\) would give
\[
 W(L'\setminus\{s\})=W(L')-w_s<W(L'),
\]
contradicting minimality.  Thus every global minimum is one-deletion
irreducible.  A rechecked safe-deletion trace is feasible because its exact
compressed-state equalities compose, but composition alone is not an
optimization argument.
\end{proof}

The converse fails in the exact \((2,2,3)\) construction.  Under identity
closure and the common zero frontier, a retained active set is safe exactly
when it covers both \(m_1\) and \(m_2\).  The singleton carriers \(s_1,s_2\)
have weights \(2,2\), and their bundle \(s_{12}\) has weight \(3\).  The five
safe active sets are
\[
 \{s_1,s_2\},\quad \{s_{12}\},\quad
 \{s_1,s_{12}\},\quad \{s_2,s_{12}\},\quad
 \{s_1,s_2,s_{12}\},
\]
with respective burdens \(4,3,5,5,7\).  Deleting the unique heaviest safe
strategy first leaves the two singleton carriers and burden four, while the
bundle-only library is the unique global minimum with burden three.  Both
endpoints preserve the source frontier and closure; the failure is global
resource optimality, not productive value.

\subsection{Sharp frontier-only loss}
\label{app:bridge-catalog}

\begin{proof}[Proof of \cref{thm:normalized-frontier-pruning-loss}.]
Use the canonical bridge catalog of Section~4: the retained library has
frontier \(F_L=(3,3)\), closure \(C_L=\{m\}\), and a dominated unique carrier
of \(m\).  Deleting that carrier preserves the frontier but makes the sole
descendant mass zero.  Conditional on retention and commitment to the unique
project, sequential generation and admission give admitted descendant mass
\(\rho^d\pi\).  Terminal reward \(R\) arrives after \(d\) dates, so the
retained one-project envelope is
\[
 \max\{0,\beta^d\rho^d\pi R-\kappa\},
\]
whereas the frontier-only comparator is zero.  The incumbent operating block
cancels under frontier equality.  At \(R=C\), the worthwhile condition selects
the research branch and yields
\[
 \operatorname{Loss}_d=\beta^d\rho^d\pi C-\kappa.
\]
For \(0\le R\le C\), monotonicity of the two-branch envelope gives the same
quantity as an upper bound, and \(R=C\) attains it.  Thus the bound is sharp
within the stated canonical comparison.  Substituting
\((d,\beta,\rho,\pi,\kappa,C)=(1,1/2,1,1,0,2M)\) gives
\cref{thm:scaled-frontier-pruning-loss}.
\end{proof}

\subsection{Capacity-constrained value}
\label{app:capacity-proof}

\begin{proof}[Proof of \cref{thm:capacity-value}.]
The inactive-only library makes every nonnegative budget feasible.  The
eligible family is finite, so each feasible subfamily has a maximizing
library.  If \(B_1\le B_2\), the first feasible family is contained in the
second, proving monotonicity.  Order the finite set of attainable burdens.
Between consecutive burdens the feasible family is unchanged, hence its
maximum is unchanged; every strict value change therefore occurs at an
attainable burden, and every exact forward difference is nonnegative.
\end{proof}

Increasing increments remain possible.  In the exact two-unit
joint-capability example,
\[
 V^\star(0)=0,\qquad V^\star(1)=0,\qquad V^\star(2)=1,
\]
so the successive unit increments are \(0\) and \(1\).  The first unit cannot
assemble the required pair, while the second completes it.  This also shows
why capacity value need not be concave; separately, the discrete price problem
may omit an attainable capacity optimum.

\subsection{Penalized finite envelope}
\label{app:penalized-proof}

\begin{proof}[Proof of \cref{thm:penalized-envelope}.]
At fixed \((\theta,b)\), each eligible library supplies the affine branch
\[
 J_{L'}(\lambda)=V_\theta(b,L')-\lambda W(L').
\]
The finite maximum is attained and is continuous, convex, and piecewise
affine, with changes of maximizing face only at finitely many pairwise
intersections.  Nonnegative burdens make every branch, and therefore the
envelope, nonincreasing.

Let \(L_1\) optimize at \(\lambda_1\) and \(L_2\) at
\(\lambda_2>\lambda_1\).  Optimality gives
\begin{align*}
 V_1-\lambda_1W_1&\ge V_2-\lambda_1W_2,\\
 V_2-\lambda_2W_2&\ge V_1-\lambda_2W_1.
\end{align*}
Adding yields \((\lambda_2-\lambda_1)(W_1-W_2)\ge0\), hence
\(W_2\le W_1\).  This applies to every optimizer pair and proves the
set-valued burden order.  Strict dominance on a neighborhood identifies one
affine branch and its slope \(-W(L')\); no derivative is selected at a tie,
and the proof imposes no nesting of raw libraries.
\end{proof}

The supporting conditional replacement formulation is in Online
Supplement~S3.
